\documentclass[11pt]{article}
\usepackage[margin=1in]{geometry}
\usepackage{amsmath,amssymb,amsthm,mathtools}
\usepackage[hidelinks]{hyperref}
\allowdisplaybreaks[2]
\setlength{\emergencystretch}{3em}
\newtheorem{theorem}{Theorem}[section]
\newtheorem{lemma}[theorem]{Lemma}
\newtheorem{proposition}[theorem]{Proposition}
\theoremstyle{remark}\newtheorem{remark}[theorem]{Remark}
\newcommand{\E}{\mathbb E}
\newcommand{\dd}{\,d}
\newcommand{\Rhalf}{\mathcal R}
\newcommand{\Lhalf}{\mathcal B}
\newcommand{\one}{\mathbf1}
\newcommand{\ones}{\boldsymbol1}
\newcommand{\eps}{\varepsilon}
\title{Moving-overlap strong passage to half mass}
\author{Analytical increment for review}
\date{Version 1; October 6, 2026}
\begin{document}\maketitle
\begin{abstract}
For the original initialized flow with $S_0=q^{10}$, a moving
overlap makes every right-half-circle displacement small.
We prove its passage to strong mass $1/2$, with a logistic clock
defect $O(q^2\log^2(1/q))$, uniform relative shape transport,
and relative radial-weight control. The fixed initial left half-circle
has mass at most $Cq^2\log^2(1/q)$ times the strong mass, so its
integrated scaled field cost is $O(q\log^2(1/q))$.
Two ingredients are proved here: a bounded-chart finite-noise
value-and-receiver-derivative reduction, and a left-ancestry
mass estimate. Neither is an additional retained hypothesis.
The chart comes from the stationary leading center, independently
of LCRC's guards. All conclusions stop at strong half mass.
\end{abstract}

\section{Model, notation and statement}
All new labels use prefix \texttt{an02momcv1:}.
Use the exact balanced, untied initialized model of P and OWPT:
$P=(P_1+P_2)/2$, $P_1=N(e_1,q^2I_2)$,
$P_2=N(\rho e_2,q^2I_2)$, $\rho\in[13/20,7/10]$.
Initial labels $\beta$ have law $\dd\beta/(2\pi)$,
$m(0,\beta)=q^{10}$ and $\theta(0,\beta)=\psi(0,\beta)=\beta$.
The fixed strong ancestry is $\Rhalf=[-\pi/2,\pi/2]$;
$\Lhalf$ is its complement in the label circle. Write
\[
 M(t)=\int_{\Rhalf}m(t,\beta)\frac{\dd\beta}{2\pi},\qquad
 \mu(t)=\int_{\Lhalf}m(t,\beta)\frac{\dd\beta}{2\pi},
 \qquad S=M+\mu .
\]
These are fixed-label masses, not moving chart masks.
Set
\[
 K(z)=\phi(z)+z\Phi(z),\quad \lambda=\rho^2,\quad
 a=\rho K(\rho a),\quad \alpha=\lambda\Phi(\rho a)/2,\quad
 d=1/2-\alpha,\quad \omega=(1-\lambda)/2,\quad s=d/\omega .
\]
The exact target-only root is $a_q=a+O(q^2)$.
Hats always mean that exact comparison flow.
For $\ell=\log(1/q)\ge1$ define
\begin{align}
 t_m&=(2/\omega+1/2)\ell,\quad C_m=\ell/2,\quad
 \sigma=10-2/\omega,\quad \zeta=\sigma-1,\nonumber\\
 u&=\sigma-\tfrac12,\qquad v=\zeta-\tfrac12=u-1,\qquad
 \eps_m=q^{s+d/2},\qquad M_0=M(t_m),\qquad
 A_q=M_0/q^u. \label{an02momcv1:scales}
\end{align}
Here $u\ge169/102$, $v\ge67/102$, $s<2$, and
$d>17503/50000$. All estimates are uniform on the ratio box.
The reference is the actual original characteristic with label
$\beta_c=qa_q$. In the strong chart write
\[
 (c(t),z(t))=(\theta(t,\beta_c)/q,\psi(t,\beta_c)/q),\quad
 h_x(t,\beta)=\theta(t,\beta)/q-c(t),\quad
 h_y(t,\beta)=\psi(t,\beta)/q-z(t).
\]
It is never reset. Define two error scales
\begin{align}
 \mathcal W_q&=q^2\ell^2,\nonumber\\
 \mathcal D_q&=q^{d/2}+\eps_m+q^v+q^2(1+\ell)+q\ell^2=o(1),
 \qquad
 \mathcal G(M_1,M_0)=
 (M_1/M_0)^{-d}\left(\frac{1-M_0}{1-M_1}\right)^\alpha .
 \label{an02momcv1:errors}
\end{align}

\begin{theorem}[Original-flow moving overlap and half-mass passage]
\label{an02momcv1:main}
For sufficiently small $q$, the following hold with no additional
entry, outside-mass, clock, or residual assumptions.

At $t_m$, OWPT's multiplicative comparisons hold with
\[
 \eta_m=O(q^v),\qquad r_m=O(q^u),\qquad
 \max_{\beta\in\Rhalf}(|h_x|+|h_y|)\le Cq^{d/2}.
\]
The normalized tail has index $3/s$, parameter $\eps_m$, and
cutoff $Cq^{d/2}$. Moreover $0<c_A\le A_q\le C_A$.

There is a first time $t_h>t_m$ with $M(t_h)=1/2$.
For every $t\in[t_m,t_h]$,
\begin{align}
 \frac{M'}M&=1-M+\delta_M,\qquad
 \int_{t_m}^{t_h}|\delta_M|\,\dd t\le C\mathcal W_q,
 \label{an02momcv1:massclaim}\\
 t_h&=10\ell-\log A_q+O(\mathcal W_q+q^u),
 \label{an02momcv1:timeclaim}\\
 |\theta(t,\beta)/q|,\ |\psi(t,\beta)/q|
 &<R:=79/51 \quad(\beta\in\Rhalf).
 \label{an02momcv1:chartclaim}
\end{align}
The reference satisfies, with a fixed $g>0$ and $\tau=t-t_m$,
\begin{align}
 |(c,z)-(a,a)|_2
 &\le C(q^2+q^v)e^{-g\tau}
   +C\eps_m\int_0^\tau e^{-g(\tau-w)}e^{-w/20}\dd w
   +Cq^2+Cq\ell^2M(t), \label{an02momcv1:meanclaim}\\
 \int_{t_m}^{t_h}|(c,z)-(a,a)|_2\dd t
 &\le C\big(q^2+q^v+\eps_m+q^2\ell+q\ell^2\big).
 \label{an02momcv1:meanintegral}
\end{align}

For every pair $\beta>\gamma$ in $\Rhalf$, for $j=x,y$,
the ratio of its positive coordinate separation at $t$ to that at
$t_m$ is
\begin{equation}\label{an02momcv1:pairclaim}
 \mathcal G(M(t),M_0)\exp[O(\mathcal D_q)],
\end{equation}
uniformly in the pair and time, including arbitrarily close pairs.
This includes signed displacements relative to $\beta_c$;
the zero displacement stays zero.
For all $\beta,\gamma\in\Rhalf$,
\begin{equation}\label{an02momcv1:weightclaim}
 \left|\log\frac{m(t,\beta)/m(t_m,\beta)}
                    {m(t,\gamma)/m(t_m,\gamma)}\right|
 \le C\mathcal W_q.
\end{equation}
The precise normalized tail sandwich is given in
\eqref{an02momcv1:tailsandwich} below. Finally,
\begin{equation}\label{an02momcv1:outsideclaim}
 \mu(t)\le Cq^2\ell^2M(t),\qquad
 \int_{t_m}^{t_h}\frac{\mu(t)}q\dd t\le Cq\ell^2.
\end{equation}
In particular the whole left ancestry, whether charted or not,
has mass $o(q)$ throughout this interval.
\end{theorem}

\section{Part A: uniformity and the moving overlap}
\begin{lemma}[Uniformity in the post-arrival time]
\label{an02momcv1:uniformity}
PROFILE's displacement, cutoff and radial comparisons have constants
independent of the terminal offset throughout their stated domain
\[
 T\ge (2/\omega)\ell+C,\qquad q^2T\le1 .
\]
For the moving time in \eqref{an02momcv1:scales}, they and OWPT
give the conclusions of Part A in Theorem~\ref{an02momcv1:main}.
\end{lemma}
\begin{proof}
Here is the needed time-uniformity audit; the completed arrival
and tail arguments are not repeated. PROFILE's local equation is
$h'=-b_q(h)h$, with $b_q(0)=\sigma_q=d+O(q^2)$,
$b_q\ge\omega/2$ and $|b_q(h)-\sigma_q|\le C|h|$ on its fixed
chart. Thus after arrival
\[
 \int_0^\infty |h(w)|\dd w\le2|h(0)|/\omega,\qquad
 h(T)=h(T_L)e^{-\sigma_q(T-T_L)}e^{O(1)}
\]
with constants independent of $T-T_L$. The root and finite chart
are uniform in $\rho$. Its Q1 arrival formula then gives
\[
 \widehat h(T,\beta)\asymp
 e^{-\sigma_q T}[q(\cos\beta+q)]^{-\sigma_q/\omega}.
\]
Replacing $\sigma_q$ by $d$ costs
$\exp\{O(q^2[T+\ell])\}$, uniformly bounded on the stated domain.
The Q4 formula uses arrival rate $1/2$ and the same local argument.
The radial proof has only the time costs
$\int_0^T d_q\le q^2T$ and
$T K(-1/(\sqrt2q))$ after the outer layer; both are uniformly
bounded for $T\le q^{-2}$. Hence the radial comparison constants
also do not depend on the offset. Normalization and the tail
change of variables preserve this uniformity.

There is no assertion that the formulas with exponent $d$ and
radial factor $e^T$ are uniform for arbitrarily large $T$ at
fixed $q$. An all-future version keeps the exact rates:
$\sigma_q$ for shape and
$\kappa_q=2g(qa_q)=1+O(q^2)$ for radial mass. Indeed
$\partial_x[2g(qx)]=4q^2 f_q(x)$, so in the local chart
$|2g(q(a_q+h))-\kappa_q|\le Cq^2h^2$.
Its time integral is bounded independently of the elapsed
post-arrival time. This also explains exactly the $q^2T$ cost
of using the leading rates. The moving logarithmic time is
strictly within PROFILE's proved domain.

At that time
\[
 \overline S(t_m)=q^u e^{2q^2t_m}=O(q^u),\qquad
 \frac1q\int_0^{t_m}\overline S(t)\dd t
 =\frac{\overline S(t_m)-q^{10}}{q(1+2q^2)}=O(q^v).
\]
OWPT's Jacobian bound and normalized-tangent proof have absolute
constants multiplying these expressions, with no separate factor
depending on the terminal offset. Reuse that proof, rather than
substitute a growing number into its fixed-offset $O$ notation.
It gives all pairwise angular comparisons with exponent $O(q^v)$,
the radial comparison with exponent $O(q^u)$, and reference
offset $|(c,z)-(a_q,a_q)|_2=O(q^v)$.

Explicitly, putting $c_0=\cos\beta$, the upper Q1 region has
\begin{equation}\label{an02momcv1:entryprofile}
 h_x,h_y\asymp \eps_m(c_0+q)^{-s},\qquad
 m(t_m,\beta)\asymp q^u(c_0+q)^2 .
\end{equation}
The weight formula holds on all $\Rhalf$. On the lower Q1 part
$|h_j|\le C\eps_m$. For Q4 the corresponding parameter is
$\eps_{-,m}=q^{(1+\lambda)s+d/2}$ with exponent
$s_-=(1-\lambda)s=2d$; its cutoff is
$Cq^{2\lambda s+d/2}$.
The Q1 cutoff is $\eps_mq^{-s}=q^{d/2}$.
These give the uniform all-label bound, the tail index $p=3/s$,
and $c_A\le M_0/q^u\le C_A$.
Let $\nu_m$ be the original normalized strong radial law. Integration
of the tail bound, using $p>3/2$, also gives the useful mean budget
\begin{equation}\label{an02momcv1:entrymoment}
 \int_{\Rhalf}(|h_x|+|h_y|)\dd\nu_m\le C\eps_m .
\end{equation}
All divisions for pair comparisons are by positive target
separations; no absolute error is divided by a small shape.
\end{proof}

\section{Left-ancestry entry mass}
\begin{lemma}[A whole-left-circle bound at overlap]
\label{an02momcv1:leftentry}
For $b_0:=\mu(t_m)/M_0$ one has
\begin{equation}\label{an02momcv1:bzero}
 0<b_0\le Cq^2\ell^2 .
\end{equation}
\end{lemma}
\begin{proof}
Use PROFILE's exact target-only Cartesian equations, valid for all
angles. Set $p_j=(\widehat U_j)_+$ and
$n_j=(-\widehat U_j)_+$. On $\Lhalf$, $p_1(0)=0$ and
$p_2(0),n_j(0)\le q^5$. Since $K(w)\le w_++\phi(0)$ and
$0\le d_q\le q^2$, their upper Dini derivatives satisfy
\[
 D^+p_1\le(1/2+q^2)p_1+Cq\sqrt{\widehat m},\qquad
 D^+p_2\le(\lambda/2+q^2)p_2+Cq\sqrt{\widehat m},\qquad
 D^+n_j\le q^2 n_j .
\]
These also hold at coordinate zero. The target spectral bound
$\sqrt{\widehat m}\le q^5e^{(1/2+q^2)t}$ yields
\[
 p_1\le Cqt\,q^5e^{(1/2+q^2)t},\quad
 p_2\le q^5e^{(\lambda/2+q^2)t}
          +Cq\,q^5e^{(1/2+q^2)t},\quad
 n_j\le q^5e^{q^2t}.
\]
The second integral uses $1-\lambda\ge51/100$.
Therefore, labelwise on the whole left circle,
\begin{equation}\label{an02momcv1:lefttarget}
 \widehat m(t,\beta)\le
 Cq^{10}e^{2q^2t}\{1+e^{\lambda t}
                     +q^2(1+t)^2e^t\}.
\end{equation}
P's early radial comparison is uniform over all initial labels,
not only $\Rhalf$, and applies up to $t_m$ since
$\overline S(t_m)=o(1)$. It changes this bound by $e^{O(q^u)}$.
Divide its integral at $t_m$ by $M_0\asymp q^{10}e^{t_m}$.
The three resulting terms are
$q^{2/\omega+1/2}$,
$q^{4+(1-\lambda)/2}$ and $Cq^2\ell^2$.
This proves \eqref{an02momcv1:bzero}. Positivity follows from
positive initial radial mass and the exact multiplicative radial
ODE on finite time intervals.
\end{proof}

\section{A finite-noise strong-chart reduction}
This section proves the reduction needed here; P's leading-order
derivation is not used as a remainder theorem.
For a measure of strong states $(\xi,\eta)$ with mass $M$ supported
in $[-R,R]^2$, write
\[
 Y=\int\eta\dd\mu_s,\quad
 F(t,r)=\phi(\min(t,r))+r\Phi(\min(t,r)),\quad
 \mathcal S(x)=\int F(\rho x,\rho\xi)\dd\mu_s .
\]
The measure and its weights are held fixed in receiver derivatives.
Define
\begin{align}
 f_x(x,y)&=\tfrac12[-(1-M)x+\rho\phi(\rho x)
                         -\rho\mathcal S(x)+\lambda y\Phi(\rho x)],
 \nonumber\\
 f_y(x,y)&=\tfrac12[-(1-M)y-Y+\rho K(\rho x)].
 \label{an02momcv1:leadingfield}
\end{align}

\begin{lemma}[Uniform value, receiver derivative and radial remainder]
\label{an02momcv1:reduction}
For any fixed finite $R$, $M\le1$, and arbitrary complementary
population of radial mass $\mu$, the exact original strong receiver
velocity $\mathcal F_q$ obeys, on $[-R,R]^2$,
\begin{align}
 \|\mathcal F_q-f\|_{C^1_{x,y}}
 &\le C_Rq^2+160\mu/q, \label{an02momcv1:c1}\\
 (\log m)'&=1-M+\delta_\beta,\qquad
 |\delta_\beta|\le C_Rq^2+3\mu . \label{an02momcv1:radial}
\end{align}
Here $160$ accommodates PROFILE's componentwise constant $40$
in the vector/matrix norms and their sum. The constants
depend only on $R$ and the ratio box.
\end{lemma}
\begin{proof}
Put $a_x=(\cos qx,\sin qx)$ and $b_y=(\cos qy,\sin qy)$;
$L_x=a_x\cdot X$. The target angular terms are
\[
 q^{-1}\E[(b_y\cdot X)(a_x^\perp\cdot X)\one_{L_x>0}],
 \qquad q^{-1}\E[(b_y^\perp\cdot X)(L_x)_+].
\]
A unit-mass strong donor $(\xi,\eta)$ subtracts
\begin{equation}\label{an02momcv1:donorkernels}
 \frac{b_y\cdot b_\eta}{q}
   \E[(L_\xi)_+(a_x^\perp\cdot X)\one_{L_x>0}],
 \qquad
 \frac{b_y^\perp\cdot b_\eta}{q}
   \E[(L_\xi)_+(L_x)_+].
\end{equation}
The target radial rate is $2\E[(b_y\cdot X)(L_x)_+]$;
the donor subtracts
$2(b_y\cdot b_\eta)\E[(L_\xi)_+(L_x)_+]$.

On $P_1$, all these gates can be removed with exponentially small
errors in value and first receiver derivative. To see the derivative
claim, differentiation creates a boundary integral on $L_x=0$.
For $|x|\le R$ and small $q$, its Gaussian mean-to-standard-deviation
ratio is at least $1/(2q)$. Conditional Gaussian moments of any
fixed degree on this boundary are bounded by a polynomial in
$q^{-1}$ times $e^{-1/(8q^2)}$. The same bound controls gate-removal
errors by Gaussian tail integration. A donor positive part is
bounded by $|X|$ and causes no worse estimate. These quantities,
even after the displayed division by $q$, are $O_R(q^2)$.
Using $\E_{P_1}XX^T=\operatorname{diag}(1+q^2,q^2)$ and retaining
the mixture factor $1/2$, the target angular terms are
$(-x,-y)/2+O_{C^1}(q^2)$, the donor terms in
\eqref{an02momcv1:donorkernels} are
$(-x,\eta-y)/2+O_{C^1}(q^2)$, and the target/donor radial
terms are $1+O(q^2)$ and $1+O(q^2)$ respectively.

For $P_2$ write $X=(qZ,V)$, $V=\rho+qZ_2$, with independent
standard normals. Let $c_x=\cos(qx)$ and $t_x=\tan(qx)/q$.
Then
\[
 L_x=qc_x(Z+t_xV),\quad
 a_x^\perp\cdot X=c_xV-q\sin(qx)Z,\quad
 b_y\cdot X=qc_y(Z+t_yV).
\]
Uniformly on the chart, $c_x=1+O_{C^1}(q^2)$,
$t_x=x+O_{C^1}(q^2)$ and $q\sin(qx)=O_{C^1}(q^2)$.
The exact one-dimensional Gaussian identities are
\[
 \E_Z[Z\one_{Z+t>0}]=\phi(t),\quad
 \E_Z[(Z+t)_+]=K(t),\quad
 \E_Z[(Z+r)_+\one_{Z+t>0}]=F(t,r).
\]
They reduce the two target angular terms to
\[
 \tfrac12\E_V[V\phi(t_xV)+t_yV^2\Phi(t_xV)],
 \qquad \tfrac12\E_V[V K(t_xV)],
\]
up to $O_{C^1}(q^2)$. All discarded Gaussian moments and their
first receiver derivatives have polynomial bounds in $|V|$,
uniformly in the compact parameters, so this error assertion
follows directly from the displayed trigonometric remainders.
Taylor expansion in $V$ at $\rho$, with
$\E(V-\rho)=0$ and $\E(V-\rho)^2=q^2$, gives
$[\rho\phi(\rho x)+\lambda y\Phi(\rho x)]/2$
and $\rho K(\rho x)/2$, with $O_{C^1}(q^2)$ errors.

For the first donor term the same calculation gives
$\tfrac12\E_V[V F(t_xV,t_\xi V)]$.
Its replacement by $\rho F(\rho x,\rho\xi)/2$ has uniform
$C^1_x$ error $O(q^2)$ even at $x=\xi$.
Here is the threshold audit. On $V\ge\rho/2$ the smaller threshold
is $V\min(t_x,t_\xi)$ and
\[
 \partial_x[V F(t_xV,t_\xi V)]
   =t_x' V^3(t_\xi-t_x)_+\phi(t_xV).
\]
This expression and its first two $V$ derivatives have uniform
polynomial bounds, including at equality of the thresholds.
The value has the same bounds. Taylor in $V$ is therefore uniform.
The event $V<\rho/2$ has exponentially small polynomial moments;
its contribution, including the receiver derivative, is negligible.
For replacing $t_x,t_\xi$ by $x,\xi$, use
$F_t(t,r)=(r-t)_+\phi(t)$ and its locally Lipschitz dependence
on $(t,r)$, together with the trigonometric bounds.
Thus no differentiation across an unbounded gate density or
unjustified second receiver derivative is used.

For the second donor angular term,
$b_y^\perp\cdot b_\eta=\sin(q(\eta-y))=O_{C^1}(q)$
and $\E_{P_2}(L_\xi)_+(L_x)_+=O_{C^1}(q^2)$.
Its divided contribution is $O_{C^1}(q^2)$.
The $P_2$ radial terms are likewise $O_R(q^2)$.
Integration over the donor measure multiplies remainders by $M\le1$.
Combining the two clusters proves the strong-only claims.
Finally PROFILE Proposition 4.1 bounds the arbitrary complement's
scaled value and derivative by $40\mu/q$ per component, and its
radial contribution by $2a_*\mu\le3\mu$, where $a_*=1/2+q^2$.
This proves both estimates.
\end{proof}

In particular the exact leading receiver Jacobian is
\begin{equation}\label{an02momcv1:jacobian}
 Df(x,y)=
 \begin{pmatrix}
 -(1-M)/2+\frac{\rho^3}{2}(y-x-\Pi(x))\phi(\rho x)
       &\frac\lambda2\Phi(\rho x)\\
 \frac\lambda2\Phi(\rho x)&-(1-M)/2
 \end{pmatrix},
 \qquad \Pi(x)=\int(\xi-x)_+\dd\mu_s.
\end{equation}
This verifies the $y$-row explicitly. Its diagonal is
$-(1-M)/2$, not $-1/2$; the latter appears only after varying the
donor mean in the collective block. At the coherent center this
is $J_{\rm sh}(M)$, with diagonal $-(1-M)/2$ and off-diagonal
$\alpha$. Hence
\[
 J_{\rm sh}(M)-r_0(M)I=
 \alpha\begin{pmatrix}-1&1\\1&-1\end{pmatrix},\qquad
 r_0(M)=-(1-M)/2+\alpha<-1/10\quad(M\le1/2).
\]

\section{Part B: a joint first-exit proof}
\subsection{Constants fixed from the leading center}
The root bound
$a\le\rho\phi(0)/(1-\lambda)<28/51$ follows from
$K(\rho a)\le\phi(0)+\rho a$ and $\phi(0)<2/5$.
Fix $R=79/51$ before choosing any bootstrap constants.
This is the only chart radius used in Lemma~\ref{an02momcv1:reduction}.
LOC's collective field $G_M$, obtained by collapsing the strong
measure at $(c,z)$, has the stationary point $(a,a)$ and
\[
 J_M=J_1-\operatorname{diag}((1-M)d,0)\preceq J_1,\qquad
 \gamma=\frac{\alpha d}{\alpha+1/2}
 \ge\gamma_*:=\frac{1183}{20800}>0.
\]
For the last bound use $\alpha\ge169/1600$, $d>7/20$ and
$\alpha+1/2<13/20$. This invokes LOC's proved mean attraction,
not a re-proof of it. Put $g=\gamma_*/2$ and
$C_0^*=724/51$, an upper bound for LOC's $4(a+3)$.

Let $C\ge1$ be a fixed enlargement of the bounded-chart constants
in Lemma~\ref{an02momcv1:reduction}, the Gaussian Lipschitz constants,
and the norm-conversion constants below. Its choice depends only
on $R$ and the ratio box. Put
\begin{equation}\label{an02momcv1:guards}
 \eps_*=\min\{1/4,\gamma_*/(4C_0^*),1/(80C)\},
 \qquad K_B=e^8 .
\end{equation}
For the actual strong family set
\[
 e(t)=|(c,z)-(a,a)|_2,\quad
 r(t)=\sup_{\beta\in\Rhalf}\max(|h_x|,|h_y|),\quad
 \overline D(t)=\int_{\Rhalf}(|h_x|+|h_y|)\dd\nu_t,\quad
 b(t)=\mu(t)/M(t).
\]
Here $\nu_t$ is the normalized strong radial measure on fixed labels.
Part A gives $e(t_m)=O(q^2+q^v)$, $r(t_m)=O(q^{d/2})$,
$\overline D(t_m)=O(\eps_m)$ and $b(t_m)=b_0=O(q^2\ell^2)$.

Stop at the first of $M=1/2$, $e=\eps_*$, $r=\eps_*$,
$b=K_Bb_0$, or the time
$t_m+T_*^0$, where
\[
 T_*^0=3\log(1/(2M_0))=O(\ell).
\]
All initial inequalities are strict for small $q$.
On this stopped interval $e+r\le1/2$, so every receiver and donor
lies strictly inside $[-R,R]^2$. No guard from LCRC is used.

\subsection{Mass, outside population and radial weights}
Average \eqref{an02momcv1:radial} over the strong radial measure:
\[
 M'/M=1-M+\delta_M,\qquad
 |\delta_M|\le Cq^2+3\mu .
\]
The outside guard makes this error $o(1)$ pointwise.
Consequently $M'/M\ge1/3$ on the stopped interval, and
\begin{equation}\label{an02momcv1:massintegral}
 \int_{t_m}^{t}M(w)\dd w\le3(M(t)-M_0)\le3/2.
\end{equation}
The exact global matrix estimates
$\|A_\theta\|\le a_*$ and $\|C_f\|\le a_*S$ give, for every
label without an angular restriction,
$(\log m)'\le(1+2q^2)(1+S)$.
Thus
\[
 (\log b)'\le (1+2q^2)(1+M+\mu)-(1-M+\delta_M)
 \le4M+Cq^2
\]
for small $q$ under $b\le K_Bb_0$.
Integrating, using \eqref{an02momcv1:massintegral} and
$q^2T_*^0=o(1)$, yields
\begin{equation}\label{an02momcv1:bimprove}
 b(t)\le e^7b_0<K_Bb_0,\quad
 \int_{t_m}^{t}\mu\le\tfrac32 K_Bb_0,\quad
 \int_{t_m}^{t}\mu/q\le\tfrac32 K_Bb_0/q .
\end{equation}
This is a generic mass comparison; it is not a capped-gain
assumption or a post-escape radial transport theorem.
In particular the outside guard improves by a fixed factor.

For every strong label the integral of the radial error is at most
$C(q^2\ell+b_0)=O(\mathcal W_q)$.
Subtract the two label equations to prove
\eqref{an02momcv1:weightclaim} on the stopped interval.
It follows also that
\begin{equation}\label{an02momcv1:normalizedweights}
 \frac{\dd\nu_t}{\dd\nu_m}=\exp[O(\mathcal W_q)]
\end{equation}
uniformly on labels. All these estimates are valid at every
stopped endpoint, not inferred backwards from a terminal cap.

\subsection{Mean attraction and shape contraction}
For a receiver on any segment between a strong label and the reference,
\eqref{an02momcv1:jacobian} gives
\begin{equation}\label{an02momcv1:jacdef}
 D\mathcal F_q(t,x,y)=J_{\rm sh}(M(t))+\mathcal E(t,x,y),\qquad
 \|\mathcal E\|_\infty
 \le C\{e+r+M\overline D+q^2+\mu/q\}.
\end{equation}
Indeed $|x-a|,|y-a|\le e+r$,
$\Pi(x)\le M(r+\overline D)$, and $\Phi$ is uniformly
Lipschitz. On the guards $\overline D\le2r$,
$M\le1/2$, so by the choice of $\eps_*$ the right side is at
most $1/20$ for small $q$, since
$\mu/q\le K_Bb_0/(2q)=o(1)$.
The base matrix $J_{\rm sh}$ is Metzler with row sums
$r_0(M)<-1/10$. The induced infinity-norm inequality for the
difference of any label and the reference is therefore
\[
 D^+\max(|h_x|,|h_y|)\le
 -\tfrac1{20}\max(|h_x|,|h_y|).
\]
This follows by the integral of the receiver Jacobian along the
joining segment and the norm inequality for $J_{\rm sh}+\mathcal E$;
no sign of the finite-$q$ off-diagonal entries is presumed.
Hence
\begin{align}
 r(t)&\le r(t_m)e^{-\tau/20},&
 \int_{t_m}^{t}r&\le20r(t_m),\nonumber\\
 \overline D(t)&\le C\eps_m e^{-\tau/20}.
 \label{an02momcv1:decay}
\end{align}
The last inequality uses \eqref{an02momcv1:normalizedweights}
and \eqref{an02momcv1:entrymoment}: each label's maximum norm
contracts, and the sum norm is bounded by twice that maximum.
It is a normalized moment estimate, not a bound replacing the
largest width by $\eps_m$.

The original reference satisfies, by LOC's collective estimate
and the finite-noise lemma,
\begin{equation}\label{an02momcv1:meanode}
 D^+e\le-\gamma_* e+C_0^*e^2
                    +C M\overline D+Cq^2+C\mu/q
 \le-ge+C M\overline D+Cq^2+C\mu/q .
\end{equation}
For clarity, the strong donor discrepancy at the reference is
bounded by $CM\overline D$: use
$|F(t,r)-F(t,\widetilde r)|\le|r-\widetilde r|$ and
$|Y-Mz|\le M\int|y-z|\dd\nu_t$.
The full complementary population is in the separately bounded
$\mu/q$ term. There is no assumed weak-family chart.

Convolution in \eqref{an02momcv1:meanode}, followed by
\eqref{an02momcv1:decay} and \eqref{an02momcv1:bimprove}, gives
\eqref{an02momcv1:meanclaim}. To obtain its last pointwise term,
use $M(w)\le M(t)e^{-(t-w)/3}$ for $w\le t$ and
$\mu(w)/q\le Cq\ell^2M(w)$.
Tonelli applied to the positive convolution kernel also gives
\eqref{an02momcv1:meanintegral}. In particular $e=o(1)$ uniformly,
improving its guard to $e<\eps_*/2$ for small $q$.
The radius guard improves by \eqref{an02momcv1:decay}.

None of the three auxiliary guards can be the first exit.
If the explicit time bound were reached first, $M'/M\ge1/3$
would force $M\ge1/2$. Thus the first exit is the half-mass hit.
The bounds also prevent finite-time breakdown before this hit;
the exact fixed-$q$ characteristic equations have bounded fields
on this bounded-mass finite interval. This proves existence,
the chart, the mean estimates and \eqref{an02momcv1:outsideclaim}.
All later statements remain restricted to this interval.

\section{The mass clock, relative tangents and the tail sandwich}
The mass defect estimate just proved is \eqref{an02momcv1:massclaim}.
Since $1-M\ge1/2$, the exact logit identity is
\[
 \frac{\dd}{\dd t}\log\frac{M}{1-M}
 =1+\frac{\delta_M}{1-M}.
\]
It implies
\begin{equation}\label{an02momcv1:duration}
 t_h-t_m=\log\frac{1-M_0}{M_0}+O(\mathcal W_q)
 =u\ell-\log A_q+\log(1-M_0)+O(\mathcal W_q).
\end{equation}
Adding $t_m$ proves \eqref{an02momcv1:timeclaim}.
Similarly, at every intermediate endpoint,
\begin{equation}\label{an02momcv1:clockconvert}
 \exp\!\left(\int_{t_m}^{t}r_0(M(w))\dd w\right)
 =\mathcal G(M(t),M_0)e^{O(\mathcal W_q)}.
\end{equation}
Indeed the time derivative of $\log\mathcal G(M(t),M_0)$ is
$r_0(M)[1+\delta_M/(1-M)]$; its additional integral is
$O(\mathcal W_q)$.

To get relative, uniform-in-pair shape control, integrate
\eqref{an02momcv1:jacdef} using the already closed budgets:
\begin{align}
 \int_{t_m}^{t_h}\sup_{\beta\in\Rhalf}\|\mathcal E(t,\beta)\|_\infty
 \dd t
 &\le C\left\{
 \underbrace{q^{d/2}}_{\text{largest width}}+
 \underbrace{\eps_m}_{\text{donor moment}}+
 \underbrace{q^2+q^v}_{\text{initial reference}}+
 \underbrace{q^2\ell}_{\text{finite noise}}+
 \underbrace{q\ell^2}_{\text{left ancestry}}\right\}.
 \label{an02momcv1:integratedtangent}
\end{align}
Let $p(t,\beta)=(\partial_\beta(\theta/q),
\partial_\beta(\psi/q))^T$ and
$\widehat p_m(\beta)=\partial_\beta(\widehat\theta(t_m,\beta)/q)>0$.
The moving-overlap extension of OWPT gives
$p(t_m,\beta)/\widehat p_m(\beta)=\ones+O(q^v)$
componentwise, uniformly in $\beta$.
Define
\[
 w(t,\beta)=
 \frac{p(t,\beta)}
 {\widehat p_m(\beta)\exp(\int_{t_m}^{t}r_0(M)\dd w)} .
\]
Holding the common population field fixed in the receiver derivative,
\begin{equation}\label{an02momcv1:normalizedtangent}
 w'=\alpha\begin{pmatrix}-1&1\\1&-1\end{pmatrix}w
       +\mathcal E(t,\beta)w .
\end{equation}
The propagator of the first matrix is nonnegative, fixes $\ones$,
and has infinity norm one. Variation of constants and Gronwall give
$\|w-\ones\|_\infty\le C(q^v+\int\|\mathcal E\|)e^{\int\|\mathcal E\|}$.
By \eqref{an02momcv1:integratedtangent}, both components are
$e^{O(\mathcal D_q)}$ and strictly positive.
This proves positivity without assuming finite-$q$ cooperativity.
Integrate the tangent over $[\gamma,\beta]$ and compare with its
initial integral. Then use \eqref{an02momcv1:clockconvert}.
This proves \eqref{an02momcv1:pairclaim}, without a lower bound
on a pair separation. Smooth dependence on a receiver's initial
label is justified by the fixed-$q$ Gaussian first-derivative
formulas above; the population is never differentiated with
respect to an individual label.

Let $D_t=|h_x(t,\beta)|+|h_y(t,\beta)|$ and
$\widehat D_m=2|\widehat\theta(t_m,\beta)/q-a_q|$.
Let $\widehat\nu_m$ be the exact target-only normalized radial law
at $t_m$. There are
$\eta=C\mathcal D_q$ and
$r=C(q^u+\mathcal W_q)$ such that, with $G=\mathcal G(M(t),M_0)$,
\begin{equation}\label{an02momcv1:tailsandwich}
 e^{-r}\widehat\nu_m\{\widehat D_m>e^\eta z/G\}
 \le\nu_t\{D_t>z\}
 \le e^r\widehat\nu_m\{\widehat D_m>e^{-\eta}z/G\}
 \quad(z>0).
\end{equation}
Indeed $D_t=G\widehat D_m e^{O(\mathcal D_q)}$ pointwise,
and the normalized density ratio
$\dd\nu_t/\dd\widehat\nu_m=e^{O(q^u+\mathcal W_q)}$
by Part A and \eqref{an02momcv1:normalizedweights}.
This proves the claimed sandwich, including the finite cutoff.
In particular the tail parameter at time $t$ is $\eps_m G$,
its index remains $3/s$, and its cutoff is $Cq^{d/2}G$,
with the same two-sided power-law range up to uniform constants.
The exponentials tending to one concern the exact normalized
comparison profile; PROFILE's power-law comparison constants
are not asserted to tend to one.

The exact handoff and half-mass scales are
\begin{align}
 \eps_mM_0^d&=A_q^d q^{10d-s},
 \label{an02momcv1:handoff}\\
 \eps_m\mathcal G(1/2,M_0)
 &=\sqrt2\,A_q^d(1-M_0)^\alpha q^{10d-s},\nonumber\\
 q^{d/2}\mathcal G(1/2,M_0)
 &=\sqrt2\,A_q^d(1-M_0)^\alpha q^{10d-2s}.
 \label{an02momcv1:terminalscales}
\end{align}
Here $s+\sigma d=10d-s$ and $d+\alpha=1/2$.
Thus the explicit moving-offset factors cancel, while $A_q$
and its clock shift remain. This completes all assertions of
Theorem~\ref{an02momcv1:main}.

\section{Scope and stopping point}
This is an exact original-flow theorem through $M=1/2$ only.
Its ingredients beyond the previously proved overlap and local
attraction results are displayed above, including the finite-noise
receiver reduction and the left-ancestry mass estimate.
No retained input is added. In particular, a small outside entry
mass was not assumed to stay small: \eqref{an02momcv1:bimprove}
proves that fact jointly with the strong passage.
No original-flow comparison is extended using a crude $S/q$ bound
after the moving overlap; only the separately bounded complement
is charged at $\mu/q$.

The order-one constants are functions of the fixed ratio box and
the analytically selected leading center and Gaussian kernels.
They do not depend on LCRC's guards. The chart is the explicit
$R=79/51$, and $A_q\in[c_A,C_A]$ is the actual overlap
normalization, not a freely selected constant.
No saturation tail, later $\int|1-\kappa_1|$ estimate,
Q2 entry contract, capped radial gain, or bridge theorem is claimed.
Work stops at the first strong half-mass time.

\begin{thebibliography}{9}\small
\bibitem{P} P, \emph{THEOREM A ANALYTICAL PROGRESS}.
\texttt{pa:earlybounds}, \texttt{pa:targetcomparison},
\texttt{pa:leading}, \texttt{pa:resident}, \texttt{pa:spectra},
\texttt{pa:order}. Its formal reduction is not used as a remainder bound.
\bibitem{PROFILE} PROFILE, \emph{AN02 strong entry profile and tail budget v1}.
Prefix \texttt{inc:AN02:profile:v1:}; \texttt{coordinate},
\texttt{root}, \texttt{arrival}, \texttt{entry}, \texttt{rightcircle},
\texttt{cost} (Proposition 4.1).
\bibitem{OWPT} OWPT, \emph{AN02 overlap window profile transfer v1}.
Prefix \texttt{an02owptv1:}; \texttt{jacobian}, \texttt{tangents},
\texttt{transfer}, \texttt{tail}.
\bibitem{SMCP} SMCP, \emph{AN02 strong M clock passage v1}.
Prefix \texttt{an02smcpv1:}; \texttt{shortobstruction},
\texttt{leadingtest}, \texttt{handoff}.
\bibitem{LOC} LOC, \emph{AN03 local rate mean tracking v1}.
Prefix \texttt{inc:AN03:local:v1:}; \texttt{field},
\texttt{meanineq} (Lemma 3.1), and its displayed collective field.
\bibitem{ROADMAP} \emph{Unified analytic theory roadmap}, v1.4,
S0 assembly requirements; S1 step 2; ledger items 12--14.
Read-only source, not a proof premise.
\end{thebibliography}
\end{document}
